\section{Exterior prolate energy and a concrete lowest-line gate}
\label{sec:spectral-prolate-exterior}

\begin{lemma}[AC-S31: a fixed-mode exterior estimate]
\label{AC-S31}
For a fixed even prolate index $n$, let $f=\phi_{n,\lambda}$
have full-line norm one, with zero extension outside the
time interval, and put $y=\mathcal Ff$ and
$\delta_n=\|(1-P_\lambda)y\|_2^2=1-\nu_n$.
For all sufficiently large $\lambda$,
\[
 |y(\lambda)|^2\leq C_n\lambda\delta_n,
 \qquad
 |y(x)|\leq C_n\lambda|y(\lambda)|/x,
 \qquad
 |y'(x)|\leq C_n\lambda^4|y(\lambda)|/x
                    \quad(x\geq\lambda).           \tag{SE.1}
\]
For its sharply cut and dilated logarithmic tail
\[
 w_n(t)=1_{t>0}\sqrt\lambda e^{t/2}y(\lambda e^t),
\]
every fixed integer $m\geq0$ and $0<s<1/2$ satisfy
\[
             \|t^mw_n\|_{H^s}
                    \leq C_{n,m,s}\lambda^{1+4s}\sqrt{\delta_n}.
                                                        \tag{SE.2}
\]
All these estimates are unconditional.
\end{lemma}
\begin{proof}
Let $\chi_n$ denote the eigenvalue of the differential
operator $PW_\lambda$ of S25, equation \texttt{wop}.
For fixed $n$,
$0\leq\chi_n\leq C_n\lambda^2$.
For the upper bound, use in the minimum principle
any fixed $(n+1)$-dimensional space of smooth functions
supported in a fixed interior interval.
On that space its form is bounded above by $\lambda^2$
times the harmonic-oscillator form.  The lower bound
is the nonnegative time-interval differential form.
The usual common prolate indexing is the one in
S25, \texttt{wop}, and Fuchs, Theorem A, equations
(1.3)--(1.6), pp.~319--320.

The Fourier continuation $y$ solves the same differential
equation globally.  To check the boundary point in this
assertion, the coefficient $\lambda^2-x^2$ vanishes at
both endpoints.  Multiplying the distributional derivative
of the zero extension by that coefficient removes its
delta terms; differentiating again introduces no term,
since the product with the interior first derivative
also vanishes at the endpoints.  Thus
$PW_\lambda f=\chi_nf$ distributionally.
Direct Fourier differentiation shows
$\mathcal F PW_\lambda=PW_\lambda\mathcal F$ on
these distributions.  Since $y$ is entire, the resulting
equation holds classically everywhere.

For $x>\lambda$ put
\[
 p=x^2-\lambda^2,\qquad V=4\pi^2\lambda^2x^2-\chi_n.
\]
For large $\lambda$, $p,V,V'>0$, and
\[
                (py')'+Vy=0.                       \tag{SE.3}
\]
With $E=p|y'|^2+V|y|^2$, differentiation gives
\[
 \left(\frac EV\right)'
       =-\left(\frac{p'}V+\frac{pV'}{V^2}\right)|y'|^2
                         \leq0.                    \tag{SE.4}
\]
The value at the endpoint is $E/V=|y(\lambda)|^2$.
In particular $|y(x)|\leq|y(\lambda)|$.
Integrating (SE.3) from the endpoint gives
\[
 py'(x)=-\int_\lambda^xV(v)y(v)dv,
 \qquad |y'(x)|\leq C_n\lambda^3|y(\lambda)|
                         \quad(\lambda\leq x\leq2\lambda).
                                                        \tag{SE.5}
\]

We next prove the sharper endpoint-to-tail comparison
in (SE.1); the short interval estimate alone would lose
two additional powers of $\lambda$.
Set $C=2\pi\lambda^2$, $x=\lambda\cosh t$, and
\[
 Y(t)=y(\lambda\cosh t),\qquad Z(t)=\sqrt{\sinh t}\,Y(t).
\]
Equation (SE.3) becomes
\[
 Z''+\Omega(t)Z=0,\qquad
 \Omega(t)=C^2\cosh^2t-\chi_n-\frac14
                                  +\frac1{4\sinh^2t}.
                                                        \tag{SE.6}
\]
Choose a fixed sufficiently small $a>0$, and set
$t_0=a/C$, $t_1=\operatorname{arcosh}2$.
Since $\lambda\cosh t_0-\lambda\leq C'a^2\lambda^{-3}$,
(SE.5) implies
$|Y(t_0)|\geq|y(\lambda)|/2$ after making $a$ small.
On $[t_0,t_1]$ the function $\Omega$ lies between
$cC^2$ and $C_aC^2$, uniformly in $\lambda$.
Moreover
\[
 \Omega'(t)=\frac{\cosh t}{2\sinh^3t}
                              (4C^2\sinh^4t-1),    \tag{SE.7}
\]
so it first decreases and then increases, with one
minimum.  The oscillator energy
$E_Z=|Z'|^2+\Omega|Z|^2$ satisfies $E_Z'=\Omega'|Z|^2$.
On its decreasing portion,
$E_Z'\geq(\Omega'/\Omega)E_Z$; on its increasing
portion $E_Z'\geq0$.
The uniform ratio of maximum to minimum of $\Omega$
therefore gives
\[
 E_Z(t_1)\geq c_aE_Z(t_0)
       \geq c_a C^2\sinh(t_0)|Y(t_0)|^2
       \geq c_n\lambda^2|y(\lambda)|^2.           \tag{SE.8}
\]

For use further outside the endpoint put $U(x)=xy(x)$.
Its equation is
\[
 U''+b(x)U'+k(x)^2U=0,\qquad
 b(x)=\frac{2\lambda^2}{x(x^2-\lambda^2)},\qquad
 k(x)^2=4\pi^2\lambda^2+
      \frac{4\pi^2\lambda^4-\chi_n-2\lambda^2/x^2}
                                    {x^2-\lambda^2}.
                                                        \tag{SE.9}
\]
For sufficiently large $\lambda$, $k^2\geq4\pi^2\lambda^2$,
and
\[
 (k^2)'=\frac{2x}{(x^2-\lambda^2)^2}
     \left(\chi_n-4\pi^2\lambda^4
                     +\frac{4\lambda^2}{x^2}
                     -\frac{2\lambda^4}{x^4}\right)<0.
                                                        \tag{SE.10}
\]
The energy $E_U=|U'|^2+k^2|U|^2$ is decreasing, since
$E_U'=-2b|U'|^2+(k^2)'|U|^2\leq0$.
At $x=2\lambda$ the two energies $E_U$ and $E_Z(t_1)$
are comparable, with constants independent of $\lambda$:
both are equivalent to
$\lambda^2|y'(2\lambda)|^2+\lambda^4|y(2\lambda)|^2$.
This follows by inserting $\sinh t_1=\sqrt3$ in
$Z'$ and using $k(2\lambda)^2\asymp\lambda^2$,
$\Omega(t_1)\asymp\lambda^4$; the mixed terms are
absorbed by these positive diagonal terms.

On $[2\lambda,3\lambda]$, $b=O(\lambda^{-1})$,
$k^2\asymp\lambda^2$, and
\[
 E_U'\geq\left(-2b+\frac{(k^2)'}{k^2}\right)E_U.
\]
Its integral gives $E_U(x)\geq c_nE_U(2\lambda)$
throughout this interval.
Choose a smooth real cutoff $\theta$, supported in
$(2\lambda,3\lambda)$, equal to one on
$[9\lambda/4,11\lambda/4]$, with $|\theta'|\leq C/\lambda$.
Multiply (SE.9) by $\theta^2\overline U$ and integrate.
After integration by parts and taking real parts,
\[
 \int\theta^2|U'|^2
  =\int\theta^2k^2|U|^2
    +\operatorname{Re}\int(\theta^2b-2\theta\theta')U'\overline U.
\]
Young's inequality absorbs half the left side and
bounds the remainder by $C_n\lambda^2\int|U|^2$.
Adding the potential term and using the lower energy
bound on the middle interval yields
\[
 \lambda E_U(2\lambda)
       \leq C_n\lambda^2\int_{2\lambda}^{3\lambda}|U(x)|^2dx.
\]
Together with (SE.8) this proves
\[
 \delta_n\geq\int_{2\lambda}^{3\lambda}|y(x)|^2dx
     \geq c_n\lambda^{-3}E_U(2\lambda)
     \geq c_n\lambda^{-1}|y(\lambda)|^2,           \tag{SE.11}
\]
which is the first assertion of (SE.1).

For the remaining exterior upper bounds, (SE.4) at
$2\lambda$ gives $|y'|\leq C_n\lambda|y(\lambda)|$,
and hence $E_U(2\lambda)\leq C_n\lambda^4|y(\lambda)|^2$.
Its decrease and $k^2\geq4\pi^2\lambda^2$ imply
\[
 |U(x)|\leq C_n\lambda|y(\lambda)|,
 \qquad |U'(x)|\leq C_n\lambda^2|y(\lambda)|
                              \quad(x\geq2\lambda).
\]
Dividing by $x$ proves (SE.1) there, with room in
the stated derivative power.  On $[\lambda,2\lambda]$
use (SE.4)--(SE.5), which give the same stated bounds.

Finally (SE.1) and (SE.11) give, for $t>0$,
\[
 |w_n(t)|\leq A e^{-t/2},\qquad
 |w_n'(t)|\leq B e^{t/2},\qquad
 A=C_n\lambda\sqrt{\delta_n},\quad
 B=C_n'\lambda^5\sqrt{\delta_n},\quad B\geq A.   \tag{SE.12}
\]
Here are explicit fractional estimates retaining
the difference between these two constants.
For $0<r<A/B$, split the translation norm of $t^mw_n$
at $R=\log(A/(Br))$.  The derivative part and the
amplitude tail, including the interval crossing zero,
are bounded by
$C_m ABr(1+|\log(A/(Br))|)^{2m}$.
For $r\geq A/B$ its squared norm is bounded by
$C_mA^2$.  Integration against $r^{-1-2s}dr$ gives
\[
 \|t^mw_n\|_{H^s}^2
              \leq C_{m,s}A^{2-2s}B^{2s}.
\]
Indeed the first integral is reduced by $r=(A/B)v$
to $\int_0^1v^{-2s}(1+|\log v|)^{2m}dv<\infty$.
This also includes the actual sharp jump at zero.
Substitution of (SE.12) proves (SE.2).
\end{proof}

\begin{corollary}[AC-S32: a sharper original-trial bound]
\label{AC-S32}
Under $A$, the original two-mode trial of AC-S30 satisfies,
for every $\varepsilon>0$,
\[
 q_L(k_\lambda)\leq C_\varepsilon
                     \lambda^{2+\varepsilon}(1-\nu_4),
 \qquad
 a_1(\log\lambda)\leq C_\varepsilon
                  \lambda^{11+\varepsilon}e^{-4\pi\lambda^2}.
                                                        \tag{SE.13}
\]
The same upper bound holds for its even part.
\end{corollary}
\begin{proof}
The two coefficients of the trial in its normalized
prolate basis are bounded.  Apply AC-S31 to each
Fourier-continued mode and use $1-\nu_0\leq1-\nu_4$.
For every arbitrarily small $\sigma>0$, the tail $w$
in (SR.11) and its first logarithmic moment obey
\[
       \|w\|_{H^\sigma}+\|tw\|_{H^\sigma}
               \leq C_\sigma\lambda^{1+4\sigma}\sqrt{\delta_\lambda}.
\]
Use the theta and derivative bounds (SR.16)--(SR.17)
with total Sobolev order less than $\sigma$, the
sharp restriction (SR.20), and the sampling inequality
(SR.22).  The explicit scalar correction (SR.18)
has norm $O(\lambda^{-1/2}\delta_\lambda)$ and is
smaller.  Equation (SR.23) therefore gives
$q_L(k_\lambda)\leq C_\sigma\lambda^{2+8\sigma}\delta_\lambda$.
Choose $8\sigma<\varepsilon$.
The nonzero $L^2$ limit, the minimum principle,
and the classical Fuchs asymptotic (SR.6) prove
the second estimate.  The only Fuchs input here
is the same scalar leakage asymptotic already
identified in AC-S30; the endpoint estimate in
AC-S31 was proved directly, without differentiating
an asymptotic or invoking an unread WKB formula.
The even-sector assertion follows as in AC-S30.
\end{proof}

\begin{corollary}[AC-S33: one sufficient second-eigenvalue estimate]
\label{AC-S33}
Assume $A$, and suppose there exist $c,\alpha>0$
such that, for all sufficiently large $\lambda$,
\[
 a_2(\log\lambda)\geq
                   c\lambda^{12+\alpha}e^{-4\pi\lambda^2}.
                                                        \tag{SE.14}
\]
Then the actual lowest Weil eigenvalue is eventually
simple and even.  If its unit eigenvector $u_\lambda$
is given the phase aligned with the even part of
$k_\lambda$, then
\[
 \int_{\mathbb R}e^{b|x|}
       \bigl|\|K\|_2u_\lambda(x)-K(x)\bigr|dx
                          \longrightarrow0
                         \quad(0\leq b<1/2).        \tag{SE.15}
\]
Thus its scalar-normalized Fourier transform converges
to $\Xi$ uniformly on every closed substrip
$|\operatorname{Im}z|\leq b<1/2$.
The estimate (SE.14) is an explicit additional analytic
gate, not a consequence proved here.
\end{corollary}
\begin{proof}
Choose $0<\varepsilon<\alpha$ in (SE.13).
Its Rayleigh quotient is $O(\lambda^{11+\varepsilon}
e^{-4\pi\lambda^2})=o(a_2)$, since its norm is bounded
below.  Hence $a_1<a_2$ eventually, proving simplicity.
Use the even part of the trial.  If the unique bottom
line were odd, this even vector would be orthogonal
to it and would have Rayleigh quotient at least $a_2$,
a contradiction.  The bottom line is therefore even.

Let $v_\lambda$ be the unit even trial, aligned with
$K$.  The elementary spectral-gap estimate of AC-S07
now gives
\[
 \|u_\lambda-v_\lambda\|_2
          \leq C\lambda^{-(1+\alpha-\varepsilon)/2}.\tag{SE.16}
\]
Both functions are supported in $[-L,L]$.
For $b>0$, Cauchy--Schwarz bounds their weighted
$L^1$ difference by $C_b\lambda^b$ times (SE.16);
for $b=0$ the factor is $\sqrt{2L}$ instead.
These errors tend to zero for every $b<1/2$.
The original trial has the weighted-$L^1$ convergence
of AC-S13, with the scalar of AC-S30; reflection
preserves that bound.  Its normalization differs
from $\|K\|_2$ by $O(\lambda^{-1/2})$, by (SR.8),
which gives another vanishing weighted error for
$b<1/2$.  Combining the three errors proves (SE.15).
The Fourier conclusion follows by taking absolute
values under the Fourier integral.  No boundary
evaluation normalization is inferred from this argument.
\end{proof}
